\documentclass{article}
\usepackage[T1]{fontenc}
\usepackage{lmodern}
\usepackage{graphicx}
\usepackage{amsmath,amssymb}
\usepackage[a4paper,margin=2cm]{geometry}
\usepackage[absolute,overlay]{textpos}
\setlength{\TPHorizModule}{1mm}
\setlength{\TPVertModule}{1mm}
\usepackage[indonesian]{babel}
\usepackage{hyperref}
\setlength{\emergencystretch}{2em}
% unit: Halaman 59

\begin{document}

% === Halaman 59 (python/native) ===

Library Crypto.Cipher 

• Pustaka Crypto.Cipher adalah package cipher di dalam Python.

https://www.pycrypto.org/api/2.6/Crypto.Cipher-module.html 

• Di dalamnya terdapat beberapa modul beberapa cipher, salah satunya RC4 

dengan nama modul Crypto.Cipher.ARC4

• Cara memanggil Pustaka Crypto.Cipher

   dengan asumsi sudah meng-intall pycryptodome dengan menggunakan perintah 

pip:

                      >> pip install pycryptodome

• Contoh program Python untuk enkripsi pesan teks dengan RC4 pada slide berikut

59

>> from Crypto.Cipher import ARC4

\end{document}
